\section{Explicit bounds for a rational seed}
\label{sec:seed-bounds}

Corollary~\ref{cor:seed} finds its seed side by testing
\(n=2,3,\ldots\) in turn, and it bounds neither the side nor the
coefficients of the seed. This section makes both explicit for a
variant of the construction with rational maps
(Sections~\ref{subsec:rational-variant}--\ref{subsec:preprocessing-bounds}),
then bounds the bit cost and the precision of recursion with any fixed
rational seed (Section~\ref{subsec:seed-recursion}) and combines the
results into one algorithm for all sizes
(Section~\ref{subsec:uniform-setup}).

\subsection{A rational variant of the construction}
\label{subsec:rational-variant}

The construction uses roots of unity in three places (Step~4 of the
proof of Theorem~\ref{thm:effective}): the rank
decompositions~\eqref{eq:convolution-rank} of convolution tensors, the
separators, and exactification. Each of them extracts one coefficient
of a polynomial, or of a Laurent polynomial, of bounded degree, and
interpolation at the integer nodes \(1,2,\ldots\) extracts the same
coefficient with rational weights. For positive integers \(n\) and
\(1\le\ell\le n\), let
\[
 \varpi_{n,\ell,\gamma}
 =[\xi^\gamma]\prod_{\substack{1\le\ell'\le n\\ \ell'\ne\ell}}
   \frac{\xi-\ell'}{\ell-\ell'}
 \qquad(0\le\gamma<n)
\]
be the coefficients of the \(\ell\)-th Lagrange basis polynomial of the
nodes \(1,\ldots,n\), where \([\xi^\gamma]\) extracts the coefficient of
\(\xi^\gamma\). Every polynomial \(f\) of degree less than \(n\), with
coefficients in a vector space over \(\Q\), satisfies
\(f(\xi)=\sum_{\ell=1}^nf(\ell)\sum_{\gamma}\varpi_{n,\ell,\gamma}\xi^\gamma\).

\begin{lemma}[Rational interpolation]\label{lem:rational-interpolation}
\leavevmode
\begin{enumerate}
\item[(a)] For positive integers \(r,b\) and \(n=r+b-1\),
\begin{equation}\label{eq:rational-convolution-rank}
 C(r,b)=\sum_{\ell=1}^{n}
 \Bigl(\sum_{\alpha=0}^{r-1}\ell^\alpha x_\alpha\Bigr)
 \Bigl(\sum_{\beta=0}^{b-1}\ell^\beta y_\beta\Bigr)
 \Bigl(\sum_{\gamma=0}^{n-1}\varpi_{n,\ell,\gamma}z_\gamma\Bigr).
\end{equation}
\item[(b)] Let \(M\ge1\), and put
\(\varpi_\ell=\ell^{2M-2}\varpi_{4M-3,\ell,2M-2}\) for
\(1\le\ell\le4M-3\). In the proof of Proposition~\ref{prop:separator},
use the source copies \(\ell=1,\ldots,4M-3\), replace the phases
\(\zeta^{\ell g'}\), \(\zeta^{\ell(e-g)}\), \(\zeta^{-\ell e'}\) by
\(\ell^{g'}\), \(\ell^{e-g}\), \(\ell^{-e'}\), and replace the scalar
\((2M-1)^{-1}\) by \(\varpi_\ell\). These rational maps give
\[
 (4M-3)\,T\degen\bigoplus_{g=0}^{M-1}T_g\,B_X(M),
\]
with the same polynomial as in that proof; its degree is at most
\((M-1)^2\), and its exponents lie in \(\{\delta^2:0\le\delta\le M-1\}\).
\item[(c)] In Lemma~\ref{lem:exactification}, the restriction
\(E\,T^{\otimes k}\restr\widetilde T_0^{\otimes k}\) is also given by
the maps \(\Psi(\ell)^{\otimes k}\), \(1\le\ell\le E\), and the scalars
\((-1)^{\ell-1}\binom E\ell\) in place of \(1/E\).
\end{enumerate}
\end{lemma}

\begin{proof}
The idea is that each identity extracts a coefficient of a polynomial
whose degree is less than the number of nodes, which interpolation at
the nodes does exactly.

(a) The product
\(P(\xi)=\bigl(\sum_\alpha x_\alpha\xi^\alpha\bigr)
\bigl(\sum_\beta y_\beta\xi^\beta\bigr)\) has degree less than \(n\),
and its coefficient of \(\xi^\gamma\) is
\(\sum_{\alpha+\beta=\gamma}x_\alpha y_\beta\). Interpolation gives this
coefficient as \(\sum_\ell P(\ell)\varpi_{n,\ell,\gamma}\). Multiplying by
\(z_\gamma\) and summing over \(\gamma\) gives
\eqref{eq:rational-convolution-rank}, by \eqref{eq:convolution}.

(b) For an integer \(\eta\) with \(|\eta|\le2M-2\), the polynomial
\(\xi^{\eta+2M-2}\) has degree between \(0\) and \(4M-4\).
Interpolating it at the \(4M-3\) nodes and comparing the coefficients
of \(\xi^{2M-2}\) gives
\begin{equation}\label{eq:rational-separator-weights}
 \sum_{\ell=1}^{4M-3}\varpi_\ell\,\ell^{\eta}
 =\begin{cases}1,&\eta=0,\\0,&0<|\eta|\le2M-2.\end{cases}
\end{equation}
After summing over the copies, the term
\(x_\alpha y_{g,\beta}z_{g,\gamma}\) becomes
\[
 \sum_{g',e,e'}\Bigl(\sum_\ell\varpi_\ell\,\ell^{e-e'+g'-g}\Bigr)
 X_{\alpha,g'}Y_{g,\beta,e}Z_{g,\gamma,e'}.
\]
Since
\(|e-e'+g'-g|\le2M-2\), the bracket is one if
\eqref{eq:separator-fourier} holds and zero otherwise, exactly as in
Step~1 of the proof of Proposition~\ref{prop:separator}. Steps~2 and~3
of that proof use only this condition, so they apply unchanged.

(c) For a polynomial \(f\) of degree less than \(E\), with coefficients
in any vector space, the \(E\)-th finite difference
\(\sum_{\ell=0}^E(-1)^\ell\binom E\ell f(\ell)\) vanishes, so
\(f(0)=\sum_{\ell=1}^E(-1)^{\ell-1}\binom E\ell f(\ell)\). In the proof of
Lemma~\ref{lem:exactification}, apply this identity to
\(f=\widetilde T(t)^{\otimes k}\), which has degree at most \(kD=E-1\),
in place of the average over the \(E\)-th roots of unity. The nodes are
nonzero, so the Laurent entries of \(\Psi\) can be evaluated there, and
the weights go on one leg of each copy.
\end{proof}

We now build the relative scheme of Proposition~\ref{prop:relative}
with \(N=k=1\), as in Theorem~\ref{thm:quantitative}, but use part~(a)
of Lemma~\ref{lem:rational-interpolation} for every rank decomposition
of a convolution tensor (in Step~2 of the proof of
Proposition~\ref{prop:relative} and for the catalyst~\eqref{eq:kappa}),
part~(b) for every separator, and part~(c) for exactification. The
separator of relation \(i\) then has \(L_i^{\Q}=4M_i-3\) source copies,
and from these copies it produces the same
polynomial~\eqref{eq:composed-relation-polynomial} as before. Hence the
proofs of Propositions~\ref{prop:catalytic} and~\ref{prop:relative}
apply with \(L\) replaced by \(L^{\Q}\), where
\begin{equation}\label{eq:rational-counts}
 L^{\Q}=\prod_{i\in\mathcal J}(L_i^{\Q})^3,\qquad
 r_*^{\Q}=\kappa_1E^j(L^{\Q})^j(2a-1)^{3Kj},
\end{equation}
while \(m_*=M^{3j}\), \(A\), \(B\), \(E\), \(\kappa_1\) and the catalyst
are unchanged. The ranks in part~(a) are those of
\eqref{eq:convolution-rank}, and every other map of the construction
has rational coefficients (Step~4 of the proof of
Theorem~\ref{thm:effective}). Thus, when \(B>A\), the
\emph{rational variant} is a relative scheme
\(r_*^{\Q}\MM_A\restr m_*\MM_B\) with rational coefficients.

Its pooled exponent differs from \eqref{eq:pooled-exponent} only in the
copy overhead. Since \(L_i^{\Q}/M_i=4-3/M_i\) lies in \([1,4)\),
\[
 0\le\log(L^{\Q}/M^3)<3I\log4.
\]
For the parameters of Theorem~\ref{thm:quantitative}, \(I\le2a^2\) and
\(\log4<2\) give
\(\log(L^{\Q}/M^3)/K<6\log4/(a\ell_a)<12/(a\ell_a)\le1\). This is the
only consequence of Step~2 of the proof of
Theorem~\ref{thm:quantitative} that Step~5 uses, and Steps~1, 3, 4
and~6 do not involve the copy count. Hence \(B>A\), and the pooled
exponent \(p\) and the source side \(A\) of the rational variant satisfy
\begin{equation}\label{eq:rational-variant-rate}
 \frac94<\tau_a\le p<\frac94+\frac{900}{\log a},\qquad
 \log A\le181a^5\ell_a^2;
\end{equation}
the lower bound holds by the argument after \eqref{eq:tau}, since all
overheads remain nonnegative. Elsewhere in the paper, including the
examples of Section~\ref{sec:costs}, the separator keeps its
\(2M_i-1\) source copies.

\subsection{An explicit seed}
\label{subsec:bounded-seed}

We now unroll the rational variant at a side given by a formula. Let
\(\eps\) be a positive rational, put \(\eps'=\min(\eps,1/2)\), and let
\(a=2^{\ell_a}\), with \(\ell_a=\lceil3600/\eps'\rceil\ge7200\), be
the index that Corollary~\ref{cor:explicit-allowance} assigns to the
allowance \(\eps'/2\). The cap \(\eps'\le1/2\) makes the pooled
exponent less than \(5/2\), which Section~\ref{subsec:preprocessing-bounds}
uses. Build the rational variant
\(r_*^{\Q}\MM_A\restr m_*\MM_B\) with the parameters of
Theorem~\ref{thm:quantitative} for this \(a\), and with any rounding
allowed by Lemma~\ref{lem:circulation-rounding}. Let
\(n_*=\lceil AB/(B-A)\rceil+1\) be its leaf threshold
(Lemma~\ref{lem:batch}), and put
\begin{equation}\label{eq:bounded-seed-parameters}
 \Xi=r_*^{\Q}(1+r_*^{\Q}m_*)\,n_*^3,\qquad
 d_0=(\Xi+1)^{\lceil2/\eps'\rceil}.
\end{equation}

\begin{theorem}[An explicit seed]\label{thm:bounded-seed}
Unrolling the pooled recursion of Lemma~\ref{lem:batch} for this scheme
at side \(d_0\) gives a rational seed of side \(d_0\) with \(R\)
products, where
\[
 R\le\Xi d_0^p<d_0^{9/4+\eps'}\le d_0^{9/4+\eps}.
\]
With the rounding rule fixed in Section~\ref{subsec:preprocessing-bounds}
and block recursion (Lemma~\ref{lem:seed-recursion}), this gives a
second proof of Theorem~\ref{thm:main}, with the seed side given by
\eqref{eq:bounded-seed-parameters} instead of a search.
\end{theorem}

\begin{proof}
The idea is that the unrolled seed has one rank-one term for each
scalar product at the last depth of the pooled recursion, and Step~3 of
the proof of Lemma~\ref{lem:batch} bounds their number. As in the proof
of Corollary~\ref{cor:explicit-allowance}, \(\log a=\ell_a\log2>\ell_a/2\)
gives \(900/\log a<1800/\ell_a\le\eps'/2\), so
\eqref{eq:rational-variant-rate} gives \(p<9/4+\eps'/2\).

Run the pooled recursion at side \(d_0\), put \(\rho=A/B\), and let
\(\ell_{\max}\) be its last depth. Step~3 of the proof of
Lemma~\ref{lem:batch} gives, also when \(\ell_{\max}=0\),
\[
 \nu_{\ell_{\max}}\le\Bigl(1+\frac{r_*^{\Q}m_*}{r_*^{\Q}-m_*}\Bigr)
 \Bigl(\frac{r_*^{\Q}}{m_*}\Bigr)^{\ell_{\max}},
 \qquad \rho^{\ell_{\max}}d_0>\rho .
\]
Since \(r_*^{\Q}/m_*=\rho^{-p}\), the second inequality gives
\((r_*^{\Q}/m_*)^{\ell_{\max}}<(r_*^{\Q}/m_*)\,d_0^p\). With
\(r_*^{\Q}-m_*\ge1\) and \(m_*\ge1\), it follows that
\(\nu_{\ell_{\max}}<r_*^{\Q}(1+r_*^{\Q}m_*)\,d_0^p\). Every product at
the last depth has side at most \(n_*\), and the classical algorithm
computes it with at most \(n_*^3\) scalar products, each of a linear
form in the left input by a linear form in the right input
(Lemma~\ref{lem:batch}). As in the proof of Corollary~\ref{cor:seed},
these products give a decomposition of \(\MM_{d_0}\) into
\(R\le\nu_{\ell_{\max}}n_*^3\le\Xi d_0^p\) rank-one tensors, with
rational coefficients because the scheme has rational coefficients.
Finally, \(d_0^{\eps'/2}\ge\Xi+1>\Xi\) and \(p<9/4+\eps'/2\) give
\(R<d_0^{9/4+\eps'}\).
\end{proof}

\subsection{The cost of computing the seed}
\label{subsec:preprocessing-bounds}

To make the seed of Theorem~\ref{thm:bounded-seed} a function of
\(\eps\), we round by the rule of Step~2 of the proof of
Theorem~\ref{thm:effective}. We count \emph{steps}: one step is an
arithmetic operation or a comparison on two rational numbers, each
stored as a reduced pair of binary integers, or one operation on an
integer index or loop counter. Let \(b_\eps\) be the total bit
length of the numerator and denominator of \(\eps\). It is not bounded
by any function of the value of \(\eps\), and it enters the bounds
through the bit lengths of the numbers computed from \(\eps\) (Step~5
of the proof below). With the
parameters of Section~\ref{subsec:bounded-seed}, put
\begin{align*}
 \mathfrak T&=\bar w+\bar z+K+I+1,\qquad
 \mathfrak F=100(j+1)\mathfrak T,\\
 \mathfrak W&=2+b_\eps+a+K+j+\bar w+\bar z+I+E+\kappa_1
       +r_*^{\Q}+m_*+A+B+\sum_{i\in\mathcal J}(M_i+L_i^{\Q}),\\
 \mathfrak H&=\mathfrak W^{\mathfrak F},\qquad
 \mathfrak Q=\mathfrak H+d_0+2,\qquad
 \mathfrak S=\mathfrak Q^{50}.
\end{align*}
These numbers are computed from \(\eps\) by the formulas above, the
occupation systems of Proposition~\ref{prop:occupations} and the
rounding rule.

\begin{theorem}[Preprocessing cost]\label{thm:construction-resources}
The dense implementation described in the proof computes the seed of
Theorem~\ref{thm:bounded-seed} from \(\eps\) in at most
\(\mathfrak S\) steps, and outputs its three rational coefficient
arrays, which have at most \(3\Xi d_0^5\) entries. Every reduced
numerator and denominator that it computes has bit length at most
\(2^{2\mathfrak S}\). With schoolbook integer arithmetic, the whole
computation, including writing the arrays, takes
\(O(2^{8\mathfrak S})\) bit operations, with an absolute implied
constant.
\end{theorem}

\begin{proof}
The idea is to fix one simple implementation and to bound every count,
dimension and exponent that it uses by \(\mathfrak W\), every array of
the assembled scheme by \(\mathfrak H\), and the arrays of the unrolled
recursion by \(\mathfrak Q^{14}\); the bit lengths then follow because
one rational operation at most doubles them, plus one bit.

\emph{Step 1: the implementation.} Store rational numbers as reduced
pairs of binary integers. Use exact Gaussian elimination for the
occupations and for basis changes, the rounding rule of
Theorem~\ref{thm:effective} for the counts, repeated multiplication for
factorials and binomial coefficients, and binary powering for integer
powers. Compute the interpolation coefficients
\(\varpi_{n,\ell,\gamma}\) and the weights of
Lemma~\ref{lem:rational-interpolation} by polynomial multiplication,
and the primitive map entries by evaluating Laurent monomials at the
integer nodes. Form tensor products, direct sums, restrictions and
compositions as dense arrays, expanding after every operation and
exploiting no sparsity.

\emph{Step 2: the primitive data.} There are \(a-1\) row systems, each
of order at most \(3a\). By the proof of
Lemma~\ref{lem:circulation-rounding}, the rounding of a row makes at
most \(\beta_{\mathrm{tot}}\) augmentations. Since
\(\beta_{\mathrm{tot}}\) is at most the sum of the fractional parts
\(\eta_e\) of that proof, it is less than the number of nonintegral
edges of the row, so each row needs at most \(J_a+1\) augmentations,
each a breadth-first search over fewer than \(J_a+3a+3\) arcs. Every
factorial or binomial argument (a split or tripling count, \(E\), or an
interpolation size \(L_i^{\Q}\) or \(r+b-1\)) is at most \(\mathfrak W\),
determinant-basis matrices have size at most \(10a^2\), and primitive
Laurent exponents are at most \(2\max_iM_i^2\). Hence generating these
data, and computing the numbers \(\mathfrak T,\ldots,\mathfrak S\),
takes fewer than \(\mathfrak W^{20}\) steps.

\emph{Step 3: dense assembly.} One catalytic use has
\(3(\bar w+\bar z)\) convolution positions and \(3\bar w\) dot-product
positions, at most \(6\mathfrak T\) primitive factors in total. Each leg
of a primitive factor, including the combined spaces on which the
determinant bases act, has dimension at most \(10a^2\le\mathfrak W^3\).
The source-copy labels of the separators, the sector and auxiliary
labels of their outputs, and the exactification label add at most
\(3I+6I+1\le10\mathfrak T\) further factors, each of dimension at most
\(\mathfrak W\). Counting all of them at once, every leg of a local
intermediate array has dimension at most
\(\mathfrak W^{18\mathfrak T+10\mathfrak T}=\mathfrak W^{28\mathfrak T}\).
The rank decompositions of the source and of the catalyst have at most
\(3\mathfrak T\) convolution positions and \(3\mathfrak T\) rank labels,
so their legs, with the matrix and copy labels, have dimension at most
\(\mathfrak W^{12\mathfrak T+3}\). During reuse the catalyst stays a
single factor (Step~3 of the proof of Proposition~\ref{prop:relative}).
The identity factors on matrix tensors add at most four powers of
\(\mathfrak W\), since \(A^2B^2\le\mathfrak W^4\), and the copy counts
add at most two more, since by \eqref{eq:one-step}, with \(L\) replaced
by \(L^{\Q}\), they are at most
\((r_*^{\Q}/\kappa_1)\,m_*\le r_*^{\Q}m_*\le\mathfrak W^2\). Hence every intermediate leg has
dimension at most
\(\mathfrak W^{28\mathfrak T+6}\le\mathfrak W^{34\mathfrak T}\).

Before its direct-sum labels are expanded, one catalytic use takes at
most \(40\mathfrak T\) dense assembly stages: two per primitive factor
for basis changes and coordinate selection (at most \(12\mathfrak T\)),
at most \(3I\) oriented separators, at most \(3I\) reindexings, at most
\(16\mathfrak T\) tensor products or direct sums, which combine the
primitive factors and the labels, and three stages for exactification
and the ledger. Expanding the labels multiplies the number of stages by
at most the number of summands, which is at most
\(\mathfrak W^{34\mathfrak T}\). The \(j\) reuse steps and the source
and catalyst decompositions therefore take at most
\(40(j+1)\mathfrak T\,\mathfrak W^{34\mathfrak T}
\le\mathfrak W^{40(j+1)\mathfrak T}\) stages. Both the dimensions and
the number of stages are below
\(\mathfrak H=\mathfrak W^{100(j+1)\mathfrak T}\), so the assembly,
including generating and copying entries, takes fewer than
\(\mathfrak H^{20}\) steps.

\emph{Step 4: unrolling.} The assembled scheme is a restriction given
by its maps, not yet a rank decomposition of \(\MM_{d_0}\), so we unroll
the pooled recursion at side \(d_0\). We have
\(n_*\le AB+2\le\mathfrak W^3\) and
\(\Xi\le\mathfrak W^{13}\le\mathfrak H\). Since
\(\nu_{\ell+1}=r_*^{\Q}\lceil\nu_\ell/m_*\rceil
\ge(r_*^{\Q}/m_*)\nu_\ell>\nu_\ell\), every depth has at most
\(\nu_{\ell_{\max}}<\Xi d_0^p\le\Xi d_0^3<\mathfrak Q^4\) products (proof
of Theorem~\ref{thm:bounded-seed}), each of side at most \(d_0\). Above
the leaf threshold the side strictly decreases, since
\(A\lceil n/B\rceil<(A/B)n+A<n\) for \(n>AB/(B-A)\), so there are at most
\(d_0\) depths. Padding adds at most \(m_*\) zero products per depth and
at most \(B\) to each side. Hence each depth has at most
\((\nu_\ell+m_*)(n_\ell+B)^2\le2\mathfrak Q^6\) coordinates per input
matrix, and all coordinate spaces at one depth, including the
intermediate source products, have dimension less than
\(\mathfrak Q^{14}\). Propagate the linear forms of the two inputs
downward, and the output forms upward, by dense matrix multiplication,
as linear forms in the \(d_0^2\) entries of each original input; each
depth costs at most a fixed multiple of \(\mathfrak Q^{42}\) steps. Each
classical product at the last depth then selects two downward forms and
a column of the upward map, which gives one rank-one term. The at most
\(d_0\) depths, the assembly of Step~3, and writing the
\(3Rd_0^2\le3\Xi d_0^5\) output entries take fewer than
\(\mathfrak Q^{50}=\mathfrak S\) steps.

\emph{Step 5: bit lengths.} Write \(\mathfrak b(\theta)\) for the
larger of the bit lengths of the reduced numerator and denominator of a
rational number \(\theta\). A sum, difference, product or defined
quotient of two rational numbers \(\theta\) and \(\theta'\) has
\(\mathfrak b\) at most
\(2\max\{\mathfrak b(\theta),\mathfrak b(\theta')\}+1\). All
initial integer literals, including \(\eps\), can be chosen with at most
\(\mathfrak Q\) bits, since \(b_\eps\le\mathfrak W\);
factorials, powers and interpolation coefficients are computed by the
counted operations. After at most \(\mathfrak S\) steps, every rational
number therefore has
\[
 \mathfrak b\le2^{\mathfrak S}(\mathfrak Q+1)-1\le2^{2\mathfrak S}.
\]
Schoolbook multiplication, division and Euclidean gcd on integers of
this length cost \(O(2^{6\mathfrak S})\) bit operations per step.
Including loop control, indexing and output, the total is
\(O(2^{8\mathfrak S})\).
\end{proof}
